% concatenated from main.tex
\documentclass[11pt]{article}
\usepackage[margin=1.15in]{geometry}
\usepackage{amsmath,amssymb,amsthm,mathtools}
\usepackage{booktabs}
\usepackage{microtype}
\usepackage[colorlinks,citecolor=blue,linkcolor=blue,urlcolor=blue]{hyperref}

\newtheorem{theorem}{Theorem}
\newtheorem{proposition}[theorem]{Proposition}
\newtheorem{lemma}[theorem]{Lemma}
\newtheorem{corollary}[theorem]{Corollary}
\newtheorem{conjecture}[theorem]{Conjecture}
\theoremstyle{definition}
\newtheorem{remark}[theorem]{Remark}
\newtheorem{problem}[theorem]{Problem}

\DeclareMathOperator{\Aut}{Aut}
\newcommand{\Lk}{\mathcal{L}}
\newcommand{\bin}[2]{\binom{#1}{#2}}
\newcommand{\Kbal}{K_{\lfloor n/2\rfloor,\lceil n/2\rceil}}

\title{The minimum of the graph likelihood}

\author{
  Simone Severini \\
 UCL
  \and
 Eric W. Weisstein \\
  Wolfram Research
  }
\date{August 12, 2026}

\begin{document}
\maketitle

\begin{abstract}
The likelihood of a finite simple undirected graph $G$ on $n$ vertices is the probability that the uniform sequential
attachment process, which at each step joins a new vertex to a uniformly random subset of
uniformly random size of the vertices already present, outputs a graph isomorphic to $G$.
Dervovic, Mocherla and Severini conjectured that the likelihood is minimised by the balanced
complete bipartite graph. We prove that, among complete bipartite graphs of a given order,
the balanced one uniquely minimises the likelihood. Exact computation shows that
it also minimises over all graphs for every order from $6$ through $14$, and that the first
counterexample occurs at $n=15$. The
blow-up of the five cycle by independent sets of size three, equivalently the circulant on
fifteen vertices with connection set $\{1,4,6\}$, has likelihood
$0.20128\ldots$ times that of $K_{7,8}$, and it is again triangle-free. We show that the failure
is not sporadic by proving that the likelihood of the balanced complete bipartite graph is
$2^{-(1/2-1/(8\ln 2)+o(1))n^2}$, whereas the minimum over all graphs of order $n$ is
$2^{-(1/2+o(1))n^2}$, so the conjectured minimiser exceeds the minimum by a factor exponential
in $n^2$. We also determine the Shannon entropy of the process to leading order, namely
$n^2/(4\ln 2)$ bits, which shows that the conjectured minimiser is in fact more likely than a
typical output of the process. The proofs rest on a vertex deletion recurrence which evaluates
the likelihood in time $O(n\,2^n)$ and which closes on the blow-ups of any fixed base graph.
\end{abstract}

\section{Introduction}

Let $G_1=K_1$ and construct a random graph $G_t$ from $G_{t-1}$ in three substeps: (1) choose $k$ uniformly from $\{0,1,\dots,t-1\}$; (2) choose a $k$-element subset $W$ of the vertex set $V(G_{t-1})$ uniformly; and (3) add a new vertex adjacent to exactly the vertices of $W$. For a finite simple undirected graph $G$ on $n$ vertices, its \emph{likelihood} is the probability that $G_n$ is isomorphic to $G$:
\[
\Lk(G)\;=\;\Pr\big[G_n\cong G\big].
\]
This invariant was introduced by Banerji, Mansour and Severini \cite{BMS14} as a graph-theoretic analogue of the infinite monkey theorem, the underlying process being the uniform case of a family of sequential attachment models studied by Janson and Severini \cite{JS13} in the context of graph limits. The likelihood is a probability distribution on the isomorphism classes of graphs of order $n$, so $\sum_G\Lk(G)=1$, where the sum runs over one representative of each isomorphism class of graphs with $n$ vertices. It quantifies how easily a graph is produced by a mechanism that uses no information about the graph built so far. Writing $\overline G$ for the complement of $G$, explicit values include $\Lk(K_n)=\Lk(\overline{K_n})=1/n!$ and $\Lk(K_{1,n-1})=n(n!)^{-2}\sum_{i=0}^{n-1}i!$, and the invariant satisfies $\Lk(G)=\Lk(\overline G)$ together with the bounds \cite[Cor.~18]{BMS14}
\begin{equation}\label{eq:BMSbounds}
\frac{1}{\lvert\Aut(G)\rvert\prod_{i=1}^{n}\bin{i-1}{\lfloor (i-1)/2\rfloor}}\;\le\;\Lk(G)\;\le\;\frac{1}{\lvert\Aut(G)\rvert}.
\end{equation}

Two extremal problems were left open in \cite{BMS14}. Weisstein \cite{MW} computed $\Lk(G)$ for every graph of order at most $10$ in 2013 and observed that for $n\le 10$, with the exceptions $n=3$ and $n=5$, the minimum is attained by $\Kbal$ and its complement. The minimum values are recorded by numerator and denominator in \href{https://oeis.org/A234234}{A234234} and \href{https://oeis.org/A234235}{A234235}, and are listed as fractions for $4\le n\le14$ in Table~\ref{tab:min} below. (The analogous maximum numerator and denominator sequences are \href{https://oeis.org/A234236}{A234236} and \href{https://oeis.org/A234237}{A234237} \cite{OEIS}.) On this numerical evidence Dervovic, Mocherla and Severini \cite[\S5]{DMS18} proposed the following assertion. Since $n=3$ and $n=5$ were already known exceptions, we understand their intended claim to apply to the nonexceptional range $n\ge6$. It is this claim that we refute below.

\begin{conjecture}[\cite{DMS18}]\label{conj:DMS}
For $n\ge6$, if $n=2p$ the minimum of $\Lk$ over graphs of order $n$ is attained by $K_{p,p}$, and if $n=2p-1$ it is attained by $K_{p-1,p}$.
\end{conjecture}

The computations below extend all four sequences with the $n=11$ and $n=12$
terms, and the two minimum sequences with the $n=13$ and $n=14$ terms. They
verify Conjecture~\ref{conj:DMS} through $n=14$. Our central result, however,
is that the conjecture is false: we show that its first failure is at $n=15$
and prove that counterexamples exist at every sufficiently large order.

The heuristic offered in \cite{DMS18} is that these graphs have a relatively large automorphism group, which by \eqref{eq:BMSbounds} pushes $\Lk$ down, and that by Mantel's theorem they are the triangle-free graphs with the largest number of edges. Our purpose is to determine how far this reasoning reaches. The results below show both its force and its limitations.

Our first results are positive. Section~\ref{sec:bip} settles the conjecture inside the family in which it is stated.

\begin{theorem}\label{thm:bip}
For every $n\ge 4$ the function $a\mapsto \Lk(K_{a,n-a})$, $0\le a\le \lfloor n/2\rfloor$, attains its unique minimum at $a=\lfloor n/2\rfloor$.
\end{theorem}

The 2013 enumeration already covered all orders through $n=10$.
Section~\ref{sec:small} adds $n=11$ and $n=12$, thereby verifying the
conjecture for $6\le n\le 12$.  Section~\ref{sec:counter} adds the remaining
cases $n=13$ and $n=14$ using a separate large-automorphism search.  We regard
a self-contained verification in exact
arithmetic as worth recording; all our computations use integer or rational
arithmetic only.

\begin{theorem}\label{thm:small}
For $4\le n\le 12$ the minimum of $\Lk$ over graphs of order $n$ is attained exactly by the two isomorphism classes $\Kbal$ and its complement, except for $n=5$, where the unique minimiser is $C_5$. The minimum values are listed in Table~\ref{tab:min}.
\end{theorem}

\begin{table}[htbp]
\centering
\begin{tabular}{@{}rll@{}}
\toprule
$n$ & $\min_G\Lk(G)$ & minimiser\\
\midrule
$4$ & $1/36$ & $K_{2,2}$\\
$5$ & $1/270$ & $C_5$\\
$6$ & $23/259200$ & $K_{3,3}$\\
$7$ & $319/54432000$ & $K_{3,4}$\\
$8$ & $319/15240960000$ & $K_{4,4}$\\
$9$ & $76441/115221657600000$ & $K_{4,5}$\\
$10$ & $76441/145179288576000000$ & $K_{5,5}$\\
$11$ & $20692621/2012184939663360000000$ & $K_{5,6}$\\
$12$ & $20692621/11155553305493667840000000$ & $K_{6,6}$\\
$13$ & $368997387881/16080060756670792571289600000000$ & $K_{6,7}$\\
$14$ & $368997387881/386307379618259120732661350400000000$ & $K_{7,7}$\\
\bottomrule
\end{tabular}
\caption{Minimum likelihoods among all simple graphs on $n$ vertices. Minimisers are unique up to complementation. Values through $n=10$ were computed in 2013 \cite{MW,OEIS}; the values for $n=11$ through $n=14$ are new here.}
\label{tab:min}
\end{table}

The two remaining smaller orders can be settled without enumerating all graphs,
by an exact enumeration of the graphs whose automorphism groups are large enough
to meet the lower bound in \eqref{eq:BMSbounds}; details are given in
Section~\ref{sec:counter}.

\begin{theorem}\label{thm:firstfailure}
For $n=13$ and $n=14$, the minimum of $\Lk$ is attained exactly by $\Kbal$
and its complement. Consequently $n=15$ is the first order $n\ge6$ for which
Conjecture~\ref{conj:DMS} fails.
\end{theorem}

Direct exhaustive search is feasible only through $n=12$, but at $n=15$ a
counterexample does occur. Write $H[\overline{K_m}]$ for the blow-up of $H$ in which every
vertex is replaced by an independent set of size $m$, adjacency being inherited.

\begin{theorem}\label{thm:counter}
Conjecture~\ref{conj:DMS} is false for $n=15$. Indeed
\[
\Lk\big(C_5[\overline{K_3}]\big)=\frac{63977511069907}{43503039787261205233506826321920000000000}
=1.4706\ldots\times 10^{-27}
\]
while
\[
\Lk(K_{7,8})=\frac{7628328998218493}{1044072954894268925604163831726080000000000}=7.3063\ldots\times 10^{-27},
\]
so
\[
\frac{\Lk(C_5[\overline{K_3}])}{\Lk(K_{7,8})}
=\frac{1535460265677768}{7628328998218493}=0.20128\ldots<1.
\]
\end{theorem}

The witness is instructive. It is a triangle-free vertex-transitive graph, so it does not contradict the first half of the heuristic of \cite{DMS18}. However, it does contradict the second half, since $C_5[\overline{K_3}]$ has $45$ edges against the $56$ of $K_{7,8}$. Blow-ups of $C_5$ are natural candidates: they remain triangle-free, nearly balanced blow-ups retain substantial symmetry, and the underlying odd cycle keeps them far from bipartite. It is precisely this family that overtakes the complete bipartite graphs.

Section~\ref{sec:asym} explains why counterexamples must eventually occur and shows that Theorem~\ref{thm:counter} is the visible edge of an asymptotic phenomenon. Automorphism-group size contributes only $O(n\log n)$ to the logarithm of the likelihood, whereas the difference between the conjectured value and the minimum is of order $n^2$. Symmetry therefore cannot make the balanced complete bipartite graph a global minimiser for large $n$.

\begin{theorem}\label{thm:asym}
As $n\to\infty$,
\[
\log_2\frac{1}{\Lk(\Kbal)}=\Big(\frac12-\frac1{8\ln 2}\Big)n^2+O(n\log n)
=(0.31966\ldots)\,n^2+O(n\log n),
\]
while
\[
\log_2\frac{1}{\min_G\Lk(G)}=\frac{n^2}{2}+O(n\log n).
\]
Consequently $\Lk(\Kbal)\big/\min_G\Lk(G)=2^{(1/(8\ln 2)+o(1))n^2}$, and Conjecture~\ref{conj:DMS} fails for every sufficiently large $n$.
\end{theorem}

Section~\ref{sec:entropy} completes the picture by computing the entropy of the process, which identifies the likelihood of a typical output and shows that the conjectured minimiser sits above it.

\begin{theorem}\label{thm:entropy}
The Shannon entropy of the distribution $\Lk$ on isomorphism classes of graphs of order $n$ equals $\frac{n^2}{4\ln 2}+O(n\log n)=(0.36067\ldots)n^2+O(n\log n)$ bits, against the $\frac{n^2}{2}+O(n\log n)$ bits of the uniform distribution on isomorphism classes.
Moreover, if $G_n$ is the output of the process, then
\[
-\frac{1}{n^2}\log_2\Lk(G_n)\;\xrightarrow{\mathrm p}\;\frac{1}{4\ln 2}.
\]
\end{theorem}

Since $\tfrac12-\tfrac{1}{8\ln 2}<\tfrac{1}{4\ln 2}$, the graph $\Kbal$ is asymptotically \emph{more} likely than a typical output of the process, which is a second way of seeing that it cannot be a minimiser. Section~\ref{sec:remarks} collects remarks on the maximum, where we show that no analogue of Theorem~\ref{thm:bip} is available, and open problems.

The tool behind all of this is a vertex deletion recurrence for $\Lk$, stated in \S\ref{sec:prelim}. It appears in this form in \cite{MW} and we claim no novelty for it; what we add is the observation that it closes on the blow-ups of any fixed base graph, which is what makes both Theorem~\ref{thm:bip} and Theorem~\ref{thm:counter} accessible.

\section{The deletion recurrence and blow-ups}\label{sec:prelim}

For a graph $G$ of order $n$, set
\[
A(G)\;=\;n!\,\lvert\Aut(G)\rvert\,\Lk(G).
\]
By \cite[Thm.~16]{BMS14}, unwinding the process along all vertex orderings,
\begin{equation}\label{eq:ordersum}
A(G)\;=\;\sum_{\sigma\in \mathfrak{S}_n}\;\prod_{i=1}^{n}\bin{i-1}{d_i(\sigma)}^{-1},
\end{equation}
where $d_i(\sigma)$ is the number of neighbours of $\sigma(i)$ among $\sigma(1),\dots,\sigma(i-1)$. In particular $A(G)=A(\overline G)$ and $A(K_n)=A(\overline{K_n})=n!$.

\begin{proposition}[\cite{MW}]\label{prop:deletion}
$A(K_1)=1$, and for every graph $G$ of order $n\ge 2$,
\begin{equation}\label{eq:deletion}
A(G)\;=\;\sum_{v\in V(G)}\frac{A(G-v)}{\bin{n-1}{\deg_G v}} .
\end{equation}
\end{proposition}

\begin{proof}
Partition the sum \eqref{eq:ordersum} according to $\sigma(n)=v$. For fixed $v$, the first $n-1$ factors run over all orderings of $V(G)\setminus\{v\}$ and reproduce $A(G-v)$, while the last factor is $\bin{n-1}{\deg_Gv}^{-1}$.
\end{proof}

Three consequences deserve to be isolated.

\begin{corollary}\label{cor:algo}
Fix a labelling of $V(G)$. Put $F(\{v\})=1$ and, for $|S|\ge 2$,
\[
F(S)=\sum_{v\in S}\frac{F(S\setminus\{v\})}{\bin{|S|-1}{\deg_{G[S]}v}}.
\]
Then
\[
\Lk(G)=\frac{F(V(G))}{n!\,\lvert\Aut(G)\rvert}.
\]
Hence $\Lk(G)$ is computable exactly in $O(n\,2^n)$ arithmetic operations and $O(2^n)$ space.
\end{corollary}

This replaces the enumeration of the $n!/\lvert\Aut(G)\rvert$ construction paths used in \cite{BMS14}. Our implementation combines graph processing in the Wolfram Language with a compiled C helper for the main exact computation. The helper evaluates an integer-scaled form of the subset recurrence modulo several $64$-bit primes, after which the exact integer is reconstructed by the Chinese remainder theorem. At the time of writing, this implementation has been used to compute exact likelihoods for every graph represented in the development version of \texttt{GraphData} \cite{GraphData} through order $29$, as well as selected larger graphs whose structure permits effective preprocessing. In particular, it returns the value in Theorem~\ref{thm:counter} for the circulant $C_{15}(1,4,6)$. Some individual graphs of order $30$ require roughly four days on current hardware and have not yet all been computed. A separate C plugin integrated with \texttt{geng} makes the exhaustive search through order $12$ feasible. The complexity of $\Lk$ remains open; see Problem~\ref{prob:complexity}.

\begin{corollary}\label{cor:deck}
$A$ is determined by the vertex-deleted deck $\{G-v: v\in V(G)\}$, counted with multiplicity. In particular, hypomorphic graphs have equal $A$, and $\Lk$ distinguishes them only through the order of their automorphism groups.
\end{corollary}

\begin{proof}
The conclusion is immediate for graphs of order at most $2$. For graphs of order at least $3$, the right-hand side of \eqref{eq:deletion} depends only on the multiset of isomorphism types $G-v$ together with the numbers $\deg_Gv=e(G)-e(G-v)$, and $e(G)$ is determined by the deck by Kelly's lemma.
\end{proof}

The third consequence is the one we shall use. Let $H$ be a graph with $V(H)=\{1,\dots,p\}$ and let $m=(m_1,\dots,m_p)\in\mathbb{Z}_{\ge0}^p$. The \emph{blow-up} $H[m]$ is obtained by replacing vertex $i$ by an independent set $V_i$ of size $m_i$ and joining $V_i$ to $V_j$ completely whenever $ij\in E(H)$. Deleting a vertex of $H[m]$ yields $H[m-e_i]$ for some $i$, so blow-ups of a fixed base form a family closed under vertex deletion and \eqref{eq:deletion} closes on it.

\begin{proposition}\label{prop:blowup}
With $n=\sum_im_i$ and $A(H[m])=1$ when $n=1$,
\begin{equation}\label{eq:blowup}
A(H[m])\;=\;\sum_{i\,:\,m_i>0}\;\frac{m_i\,A(H[m-e_i])}{\bin{n-1}{\sum_{j\sim i}m_j}} .
\end{equation}
\end{proposition}

\begin{proof}
Every vertex of $V_i$ has degree $\sum_{j\sim i}m_j$ and deleting it gives $H[m-e_i]$; apply \eqref{eq:deletion} and group the $n$ terms into the $p$ classes $V_1,\dots,V_p$.
\end{proof}

Taking $H=K_p$ gives all complete multipartite graphs, indexed by partitions of $n$; taking $H=K_2$ gives all complete bipartite graphs; taking $H=C_5$ gives the family of Theorem~\ref{thm:counter}. For the automorphism groups of blow-ups we use the theorem of Sabidussi \cite{Sab59,Sab61} on lexicographic products, in the form given by Grech and Kisielewicz \cite[\S1]{GK22}. If no two vertices of $H$ have the same neighbourhood and no two have the same closed neighbourhood, then $\Aut(H[\overline{K_m}])=\Aut(H)\ltimes \mathfrak{S}_m^{\,p}$, of order $\lvert\Aut(H)\rvert\,(m!)^p$.

\section{Complete bipartite graphs}\label{sec:bip}

Specialising \eqref{eq:blowup} to $H=K_2$ and normalising, put
\[
b^{(n)}_a\;=\;\frac{A(K_{a,n-a})}{a!\,(n-a)!},\qquad 0\le a\le n,
\]
so that, since $\lvert\Aut(K_{a,n-a})\rvert=a!\,(n-a)!\,(1+[\,2a=n\,])$,
\begin{equation}\label{eq:Lb}
\Lk(K_{a,n-a})=\frac{b^{(n)}_a}{n!\,\big(1+[\,2a=n\,]\big)} .
\end{equation}
Proposition~\ref{prop:blowup} becomes, using $\bin{n-1}{n-a-1}=\bin{n-1}{a}$,
\begin{equation}\label{eq:brec}
b^{(n)}_a=\frac{b^{(n-1)}_{a-1}}{\bin{n-1}{a-1}}+\frac{b^{(n-1)}_{a}}{\bin{n-1}{a}}
\qquad (1\le a\le n-1),
\end{equation}
with $b^{(n)}_0=b^{(n)}_n=1$ for all $n\ge1$ and $b^{(1)}_0=b^{(1)}_1=1$. Note $b^{(n)}_a=b^{(n)}_{n-a}$.

\begin{lemma}\label{lem:b1}
For $n\ge 2$ we have $b^{(n)}_1=b^{(n)}_{n-1}=\dfrac{1}{(n-1)!}\displaystyle\sum_{k=0}^{n-1}k!$, and $1<b^{(n)}_1\le 2$ with equality if and only if $n\in\{2,3\}$.
\end{lemma}

\begin{proof}
$K_{1,n-1}$ is the star, and $\Lk(K_{1,n-1})=n(n!)^{-2}\sum_{k=0}^{n-1}k!$ by \cite[Prop.~20]{BMS14}; substituting into \eqref{eq:Lb} with $\lvert\Aut\rvert=(n-1)!$ gives the formula. For the bounds, isolate the top term and observe that
$\sum_{k=0}^{n-2}k!\le (n-2)!+\sum_{k=0}^{n-3}k!\le (n-2)!+(n-2)(n-3)!=2(n-2)!$,
which is smaller than $(n-1)!=(n-1)(n-2)!$ as soon as $n\ge4$. Hence $b^{(n)}_1=1+\sum_{k=0}^{n-2}k!/(n-1)!<2$ for $n\ge4$, while $b^{(2)}_1=b^{(3)}_1=2$; and $b^{(n)}_1>1$ because the term $k=n-1$ alone contributes $1$.
\end{proof}

\begin{lemma}\label{lem:bbound}
For all $n\ge1$ and all $a$ we have $b^{(n)}_a\le 2$. Moreover $b^{(n)}_a\le 4/(n-1)$ whenever $n\ge4$ and $2\le a\le n-2$.
\end{lemma}

\begin{proof}
Induction on $n$. For $n\le 3$ every value is $1$ or $2$ by Lemma~\ref{lem:b1}. Let $n\ge4$ and assume the first claim for $n-1$. If $a\in\{0,n\}$ then $b^{(n)}_a=1$, and if $a\in\{1,n-1\}$ then $b^{(n)}_a\le2$ by Lemma~\ref{lem:b1}. If $2\le a\le n-2$ then $1\le a-1$ and $a\le n-2$, so $\bin{n-1}{a-1}\ge n-1$ and $\bin{n-1}{a}\ge n-1$, and \eqref{eq:brec} gives $b^{(n)}_a\le 2/(n-1)+2/(n-1)=4/(n-1)\le 2$.
\end{proof}

\begin{lemma}\label{lem:diff}
For $1\le a\le n-2$,
\[
b^{(n)}_a-b^{(n)}_{a+1}=\frac{b^{(n-1)}_{a-1}}{\bin{n-1}{a-1}}-\frac{b^{(n-1)}_{a+1}}{\bin{n-1}{a+1}} .
\]
\end{lemma}

\begin{proof}
Subtract the instance $a+1$ of \eqref{eq:brec} from the instance $a$; the two middle terms $b^{(n-1)}_{a}/\bin{n-1}{a}$ cancel.
\end{proof}

\begin{proposition}\label{prop:mono}
For every $n\ge 4$ the sequence $b^{(n)}_1,b^{(n)}_2,\dots,b^{(n)}_{\lfloor n/2\rfloor}$ is strictly decreasing.
\end{proposition}

\begin{proof}
Induction on $n$. For $n=4$ we have $b_1=5/3>b_2=4/3$, and for $n=5$, $b_1=17/12>b_2=23/36$. Let $n\ge 6$, assume the statement for $n-1$, and fix $a$ with $1\le a\le \lfloor n/2\rfloor-1$. By Lemma~\ref{lem:diff} it suffices to prove
\begin{equation}\label{eq:target}
\frac{b^{(n-1)}_{a-1}}{\bin{n-1}{a-1}}>\frac{b^{(n-1)}_{a+1}}{\bin{n-1}{a+1}} .
\end{equation}

Assume first $a\ge2$, so that $1\le a-1<a+1\le\lfloor n/2\rfloor$. If $a+1\le\lfloor (n-1)/2\rfloor$ then the inductive hypothesis gives $b^{(n-1)}_{a-1}>b^{(n-1)}_{a+1}$, while $a+1\le\lfloor (n-1)/2\rfloor$ forces $\bin{n-1}{a-1}<\bin{n-1}{a+1}$, and \eqref{eq:target} follows. Otherwise $a+1>\lfloor (n-1)/2\rfloor$ together with $a+1\le\lfloor n/2\rfloor$ forces $n$ even and $a+1=n/2$, so $n-1=2\mu+1$ with $\mu=n/2-1=a$. Then $b^{(n-1)}_{a+1}=b^{(n-1)}_{n-1-(a+1)}=b^{(n-1)}_{a}$ and $\bin{n-1}{a+1}=\bin{2\mu+1}{\mu+1}=\bin{2\mu+1}{\mu}=\bin{n-1}{a}$, so \eqref{eq:target} reduces to $b^{(n-1)}_{a-1}/\bin{n-1}{a-1}>b^{(n-1)}_{a}/\bin{n-1}{a}$, which holds because $1\le a-1<a\le \lfloor (n-1)/2\rfloor$ gives $b^{(n-1)}_{a-1}>b^{(n-1)}_{a}$ by induction and $\bin{n-1}{a-1}<\bin{n-1}{a}$.

Now let $a=1$. Then $b^{(n-1)}_{0}=1$ and $\bin{n-1}{0}=1$, so \eqref{eq:target} reads $b^{(n-1)}_{2}<\bin{n-1}{2}$. Since $n\ge6$, Lemma~\ref{lem:bbound} gives $b^{(n-1)}_2\le 4/(n-2)\le1<\bin{n-1}{2}$.
\end{proof}

\begin{proof}[Proof of Theorem~\ref{thm:bip}]
By \eqref{eq:Lb} and Proposition~\ref{prop:mono}, $\Lk(K_{a,n-a})$ is strictly decreasing for $1\le a\le\lfloor n/2\rfloor-1$, and the passage from $a=\lfloor n/2\rfloor -1$ to $a=\lfloor n/2\rfloor$ decreases it as well, strictly, since for $n$ even the additional factor $(1+[2a=n])^{-1}=1/2$ only helps. It remains to compare with $a=0$, where $\Lk(\overline{K_n})=1/n!$. For $n\ge 6$ we have $2\le\lfloor n/2\rfloor\le n-2$, so Lemma~\ref{lem:bbound} gives $b^{(n)}_{\lfloor n/2\rfloor}\le 4/(n-1)<1$; for $n=4$, $b_2^{(4)}/2=2/3<1$ and for $n=5$, $b^{(5)}_2=23/36<1$.
\end{proof}

\begin{remark}
The function $a\mapsto\Lk(K_{a,n-a})$ is not monotone on $0\le a\le\lfloor n/2\rfloor$; the star $a=1$ is the maximum of the family and lies above the empty graph $a=0$. For $n=10$ the values are $2.76\times 10^{-7}$, $3.11\times 10^{-7}$, $3.63\times 10^{-8}$, $1.20\times 10^{-9}$, $2.29\times 10^{-11}$, $5.27\times 10^{-13}$.
\end{remark}

\section{Exhaustive verification for \texorpdfstring{$n\le 12$}{n <= 12}}\label{sec:small}

The structured searches below use the following elementary screening, which
combines the deletion recurrence, the lower bound in \eqref{eq:BMSbounds},
and the fact that automorphisms preserve degrees. Write
$P_n=\prod_{i=1}^{n}\bin{i-1}{\lfloor (i-1)/2\rfloor}$ and let $m_d(G)$ be
the number of vertices of $G$ of degree $d$.

\begin{lemma}\label{lem:screen}
If $n\ge2$ and $\Lk(G)\le\lambda$, then
\[
\prod_{d\ge0}m_d(G)!\;\ge\;\lvert\Aut(G)\rvert\;\ge\;
\frac{1}{n\lambda P_{n-1}}
\sum_{v\in V(G)}\bin{n-1}{\deg_Gv}^{-1}
\;\ge\;\frac{1}{\lambda P_n}.
\]
\end{lemma}

\begin{proof}
The left inequality holds because $\Aut(G)$ embeds in the direct product of
the symmetric groups on the degree classes.  Applying \eqref{eq:BMSbounds}
to each graph $G-v$ gives
\[
A(G-v)=(n-1)!\lvert\Aut(G-v)\rvert\Lk(G-v)\ge\frac{(n-1)!}{P_{n-1}}.
\]
Substitution in \eqref{eq:deletion}, followed by
$\Lk(G)=A(G)/(n!\lvert\Aut(G)\rvert)\le\lambda$, gives the middle inequality.  The
last follows from
$\bin{n-1}{\deg_Gv}\le\bin{n-1}{\lfloor(n-1)/2\rfloor}$ and the definition
of $P_n$.
\end{proof}

For the exhaustive proof we use one implementation uniformly through order
$12$.  Setting
$D_i=\operatorname{lcm}\{\bin{i-1}{k}:0\le k\le i-1\}$ and
$\widehat F(S)=F(S)\prod_{i\le |S|}D_i$, the recurrence of
Corollary~\ref{cor:algo} becomes
\[
\widehat F(S)=\sum_{v\in S}\widehat F(S\setminus\{v\})
\frac{D_{|S|}}{\displaystyle\bin{|S|-1}{\deg_{G[S]}v}},
\]
in which every factor is a positive integer.  For $n\le12$,
$\widehat F(V)\le n!\prod_{i\le n}D_i<2^{90}$.  The integrated
\texttt{geng} plugin uses checked unsigned $128$-bit arithmetic for this
recurrence, exact \texttt{nauty} automorphism orders, and exact $192$-bit
cross-products for likelihood comparisons.  It rejects larger orders.  It
enumerates every unlabelled graph through order $11$ and, using complement
symmetry, one representative from every complement orbit at order $12$.  As
a global check, it reproduces the extrema previously known through order $10$.

At order $11$, the program streams all $1\,018\,997\,864$ unlabelled graphs
and obtains
\[
\min_G\Lk(G)=\frac{20692621}{2012184939663360000000}.
\]
The two minimisers are $K_{5,6}$ and its complement. Graph6 representations
of the complement and $K_{5,6}$ are \verb|J~{?GKF@wN_| and
\verb|J??F~z{~Fw?|, respectively.

For $n=12$, complement symmetry reduces the search to one representative
from each complement orbit. There are exactly
$82\,545\,586\,656=(165\,091\,172\,592+720)/2$ such orbits, since there are
$165\,091\,172\,592$ unlabelled graphs of order $12$, of which $720$ are
self-complementary. The search obtains
\[
\min_G\Lk(G)=\frac{20692621}{11155553305493667840000000}.
\]
The two minimisers are $K_{6,6}$ and its complement, with graph6
representations \verb+KsaCB|}^b{No+ and \verb+K~~w?CB?wF_^+. This proves
Theorem~\ref{thm:small}.

\begin{remark}
The minimum is well separated. For $n=10$ the second smallest value is $8.39\times 10^{-12}$, larger than the minimum by a factor of about $16$.
\end{remark}

\section{A counterexample}\label{sec:counter}

\begin{proof}[Proof of Theorem~\ref{thm:counter}]
Let $H=C_5$ with $V(H)=\mathbb{Z}_5$ and $m=(3,3,3,3,3)$, so $G=C_5[\overline{K_3}]$ is a $6$-regular graph of order $15$ with $45$ edges; it is isomorphic to the circulant $C_{15}(1,4,6)$, the parts being the cosets of the subgroup of order $3$ in $\mathbb{Z}_{15}$. No two vertices of $C_5$ share a neighbourhood or a closed neighbourhood, so $\lvert\Aut(G)\rvert=\lvert\Aut(C_5)\rvert\,(3!)^5=10\times 6^5=77760$ by the theorem of Sabidussi. Iterating \eqref{eq:blowup} over the $4^5$ vectors $m'\le m$ in exact rational arithmetic gives
\[
A(G)=\frac{63977511069907}{427822618422142944000000},
\]
and $\Lk(G)=A(G)/(15!\times 77760)$ is the stated value. Independently, Corollary~\ref{cor:algo} applied to $G$ with $2^{15}$ subsets and rational arithmetic returns the same $A(G)$. Theorem~\ref{thm:bip} or a direct application of \eqref{eq:brec} gives $\Lk(K_{7,8})$, and the comparison is a rational inequality.
\end{proof}

Further exact computations locate the example.

\begin{proposition}\label{prop:first15}
For $n\in\{13,14\}$, no independent blow-up $H[m]$ with
$\lvert V(H)\rvert\le12$ has likelihood smaller than $\Kbal$.  For
$n\in\{16,18\}$ the same holds with $\lvert V(H)\rvert\le8$.  For $n=15$ the blow-up
$C_5[\overline{K_3}]$ does.  Also, no vertex-transitive graph
of order between $6$ and $14$ has likelihood smaller than $\Kbal$, whereas the
circulant $C_{15}(1,4,6)$ does.
\end{proposition}

For the new cases this is a finite exact computation using
\eqref{eq:blowup}.  Through base order $9$, every unlabelled base and every
positive part vector is checked.  At base order $9$ this gives totals of
$135\,960\,660$ and $353\,497\,716$ blow-ups for $n=13$ and $n=14$,
respectively.  In each case the minimum is attained by $\Kbal$.  Blow-ups of
smaller bases are included, since a part can be split by duplicating its base
vertex as a false twin.

Base order $10$ is handled by an exact automorphism screen.  A base with
false twins can be reduced by merging them and is therefore already covered
by a smaller base.  If $H$ is false-twin-free, its fibres are precisely the
false-twin classes of $H[m]$, and hence
\[
\lvert\Aut(H[m])\rvert=\left(\prod_i m_i!\right)
\lvert\operatorname{Stab}_{\Aut(H)}(m)\rvert
\le \left(\prod_i m_i!\right)\lvert\Aut(H)\rvert.
\]
There are $12\,005\,168$ unlabelled bases of order $10$, of which
$1\,990\,286$ have false twins.  A strict competitor at order $13$ must have
automorphism order at least $183$.  Since the largest fibre-factorial product
among the $220$ positive compositions of $13$ into $10$ parts is $4!=24$,
only false-twin-free bases with automorphism order at least $8$ need be kept;
there are $50\,714$.  The part-vector-specific bound leaves $2\,368\,100$
targets for an exact automorphism computation, of which $528\,765$ meet the
necessary target threshold and require likelihood evaluation.  None beats
$K_{6,7}$.  At order $14$ the corresponding target threshold is $2552$ and
the maximum fibre-factorial product is $5!=120$.  There are $7\,498$
false-twin-free bases with automorphism order at least $22$; the refined
screen computes exact automorphism orders for $526\,705$ targets and exact
likelihoods for the surviving $80\,258$.  None beats $K_{7,7}$.

At base order $11$ there are $1\,018\,997\,864$ unlabelled bases, of
which $106\,088\,988$ have false twins.  At order $13$ there are only $66$
positive part vectors and their largest fibre-factorial product is $3!=6$.
The target automorphism threshold $183$ therefore leaves only
false-twin-free bases with automorphism order at least $31$; exactly
$41\,737$ bases survive.  The part-vector-specific screen requires
$2\,210\,912$ exact target-automorphism computations and $555\,612$
likelihood evaluations.  None beats $K_{6,7}$.  At order $14$ there are
$286$ positive part vectors, the largest fibre-factorial product is $4!=24$,
and the target threshold $2552$ requires base automorphism order at least
$107$.  Exactly $5\,700$ bases survive; $497\,915$ target automorphism orders
and $79\,908$ likelihoods are evaluated, and none beats $K_{7,7}$.  The
$41\,737$ bases retained during the order-$13$ stream were saved and reused
for the order-$14$ search, so the full order-$11$ \texttt{geng} stream was
generated only once.

At base order $12$ there are $165\,091\,172\,592$ unlabelled bases, of
which $10\,454\,883\,132$ have false twins.  At order $13$ there are $12$
positive part vectors, each with fibre-factorial product $2$.  The target
threshold $183$ therefore requires base automorphism order at least $92$,
and an exhaustive $4\,096$-shard stream retains exactly $140\,474$ bases.
The resulting $1\,685\,688$ target representations require exact
automorphism computations; $927\,658$ are discarded by the target threshold
and the remaining $758\,030$ receive exact likelihood evaluations.  None
beats $K_{6,7}$.

At order $14$ there are $78$ positive part vectors and the largest
fibre-factorial product is $3!=6$, so a possible competitor must have base
automorphism order at least $426$.  The bases saved during the order-$13$
stream therefore contain every possibility; the stronger base bound retains
$9\,224$ of them.  Of the resulting $719\,472$ target representations,
$487\,416$ require exact target-automorphism computations and $121\,413$
survive for exact likelihood evaluation.  None beats $K_{7,7}$.

The corresponding searches through base order $7$ at $n=16$ and $n=18$
evaluate $5\,225\,220$ and $12\,920\,544$ blow-ups, respectively, and again
return $\Kbal$.
At $n=16$, extending the search through base order $8$ checks every positive
eight-part vector for all $12\,346$ unlabelled bases.  After omitting bases
with false twins, which reduce to smaller bases, and applying the exact
automorphism screen, it evaluates $34\,372\,082$ likelihoods; none beats
$K_{8,8}$.  The analogous base-order-$8$ search at $n=18$ evaluates all
$156\,498\,056$ likelihoods; none beats $K_{9,9}$.
As a further check at base order $9$, the minimum among all $274\,668$ graphs
$H[\overline{K_2}]$ with $\lvert V(H)\rvert=9$ is attained when $H$ is the
$3\times3$ rook graph.  It
comes close, but remains above the conjectured minimum by the exact factor
\[
\frac{4427743161718052172975}{4400906902589983024166}
=1.0060978\ldots.
\]
A separate search considers every graph $rH+tK_1$, and its complement, where
$H$ is connected of order at most $7$, $r\ge2$, and the total order is $13$
or $14$.  There are $156$ such representations at order $13$ and $1010$ at
order $14$.  At order $13$ the smallest likelihood in this family occurs at
$2K_6+K_1$ and is $17.1364\ldots$ times $\Lk(K_{6,7})$; at order $14$ the
minimum occurs at $2K_7=\overline{K_{7,7}}$.

The vertex-transitive assertion is automatic from Theorem~\ref{thm:small}
through order $12$.  At orders $13$ and $14$, respectively, enumeration from
the transitive permutation groups gives $14$ and $56$ isomorphism classes of
vertex-transitive graphs.  At order $13$ the smallest likelihood among them
is attained by the Paley graph of order $13$ and exceeds $\Lk(K_{6,7})$ by
the exact factor
\[
\frac{7061981141760}{368997387881}=19.1382\ldots.
\]
At order $14$ the only minima in this class are $K_{7,7}$ and its complement.

The lower bound in \eqref{eq:BMSbounds} also gives a useful restriction on
any unrestricted counterexample.  For $n=13,14,16,18$, respectively, a graph
$G$ satisfying $\Lk(G)<\Lk(\Kbal)$ must have
\[
\lvert\Aut(G)\rvert\ge183,\quad2552,\quad1556,\quad462.
\]
Indeed, before taking the next integer, the four lower bounds are
$182.255\ldots$, $2551.577\ldots$, $1555.153\ldots$ and $461.923\ldots$.
This makes enumeration by large automorphism group a possible alternative to
enumerating every graph.  The exact blow-up searches also rule out an
additional tie with $\Kbal$.  They show that any remaining graph of order
$13$ with likelihood at most $\Lk(K_{6,7})$, as well as its complement, must
be false-twin-free; equivalently, it can have neither false twins nor true
twins.  At order $14$, any remaining graph with likelihood at most
$\Lk(K_{7,7})$ can have at most one false-twin pair and at most one true-twin
pair.  If either occurs, the corresponding quotient of order $13$ must have
automorphism order at least $1276$.

We complete the search by enumerating permutation groups rather than graphs.
For a subgroup $\Gamma\le\mathfrak S_n$, every $\Gamma$-invariant graph is a union of
orbits of $\Gamma$ on the unordered vertex pairs.  Moreover, if two vertices $u,v$
have the property that $\{u,w\}$ and $\{v,w\}$ belong to the same pair orbit
for every $w\notin\{u,v\}$, then every such graph makes $u,v$ either false
twins or true twins, according as $uv$ is absent or present.

At order $13$, the complete table of marks for $\mathfrak S_{13}$ in
\textsc{TomLib} \cite{TomLib} contains $20\,832$ conjugacy classes of
subgroups, of which $9634$ have order at least $183$.  The forced-twin test
leaves $26$ classes and $3296$ orbital graph representations.  Direct twin
testing leaves $464$, and canonical labelling reduces these to $100$
isomorphism classes.  Their exact likelihoods have minimum ratio
\[
\frac{\Lk(G)}{\Lk(K_{6,7})}
=\frac{38307963944498}{368997387881}>1.
\]

For a graph of order $14$ with a false-twin pair, uniqueness of that pair
forces its automorphism group to preserve the pair.  After factoring out its
transposition, the induced group on the remaining $12$ vertices has order at
least $1276$.  The complete table for $\mathfrak S_{12}$ has $10\,723$
subgroup classes, of which $956$ meet this bound.  Their pair orbits, together
with their vertex orbits describing the common neighbourhood of the twin
pair, give $7\,523\,958$ representations.  Exact twin testing leaves
$6\,088\,090$, and canonical labelling leaves $122\,914$ graphs.  Their
minimum exact ratio is
\[
\frac{\Lk(G)}{\Lk(K_{7,7})}
=\frac{17952366553322}{368997387881}>1.
\]
Complementation covers the case of a true-twin pair.

It remains to treat graphs of order $14$ having neither kind of twin.  Starting
from $\mathfrak S_{14}$, we recursively enumerate conjugacy-class
representatives of maximal subgroups whose orders remain at least $2552$,
identifying subgroups conjugate in $\mathfrak S_{14}$.  Every subgroup above
the threshold occurs in such a chain.  GAP returns no failed maximal-subgroup
computations and gives $6280$ classes.  The forced-twin test leaves four
classes and $40$ orbital representations; direct twin testing and canonical
labelling leave four graphs.  Their minimum exact ratio is again
\[
\frac{38307963944498}{368997387881}>1.
\]
Together with the blow-up search, these computations prove
Theorem~\ref{thm:firstfailure}.

The same subgroup method substantially narrows the order-$16$ problem but
does not yet settle it.  Here the automorphism threshold is $1556$.  A
checkpointed maximal-subgroup traversal of $\mathfrak S_{16}$ finds
$116\,597$ subgroup classes above the threshold, with no failed
maximal-subgroup computations.  Of these, $116\,553$ force a false- or
true-twin pair in every invariant graph.  The remaining $44$ classes give
$4580$ orbital graph representations; direct twin testing leaves $738$, and
canonical labelling leaves $128$ isomorphism classes.  None has likelihood
below $\Lk(K_{8,8})$.  Thus any counterexample at order $16$ must contain
false or true twins.  The existing blow-up searches cover such graphs only
through base order $8$, so the forced-twin cases with larger quotients remain
to be checked and the unrestricted order-$16$ problem remains open.

At the degree-sequence level, the final inequality of Lemma~\ref{lem:screen}
leaves $608\,239$ of the $836\,315$ graphical sequences at order $13$ and
$1\,041\,046$ of the $3\,166\,852$ sequences at order $14$.  The strengthened
middle inequality reduces these counts to $152\,691$ and $175\,878$.
Applying the deletion recurrence once more and minimizing over every possible
neighbour-degree assignment, a relaxation that remains a rigorous lower
bound, leaves $69\,797$ and $75\,888$ sequences, or $34\,949$ and $37\,944$
orbits under complementation.  Thus degree information alone substantially
narrows the search but does not settle either order.

Thus $n=15$ is the first failure of the conjecture. Second, the advantage of
the $C_5$ blow-ups grows very quickly with $n$; for $n=25$,
\[
\Lk\big(C_5[\overline{K_5}]\big)=1.202\ldots\times 10^{-77}
\qquad\text{against}\qquad
\Lk(K_{12,13})=6.960\ldots\times 10^{-72},
\]
a factor of about $5.8\times 10^{5}$.  Exact enumeration of every positive
part vector, modulo the dihedral symmetries of $C_5$, shows that through
$n=50$ the minimum among positive $C_5$ part vectors is always attained by a
nearly balanced vector, meaning one whose coordinates differ by at most one.
It is smaller than $\Kbal$ at $n=15$, at $n=17$, and at
every order from $19$ through $50$.  Exact evaluation of the best nearly
balanced vectors continues this sequence of counterexamples without a gap
through $n=100$; selected early examples are listed in
Table~\ref{tab:c5more}.

\begin{table}[htbp]
\centering
\begin{tabular}{@{}rcr@{.}l@{}}
\toprule
$n$ & $m$ & \multicolumn{2}{c}{$\Lk(C_5[m])/\Lk(\Kbal)$}\\
\midrule
$17$ & $(3,3,3,4,4)$ & $0$ & $32136\ldots$\\
$19$ & $(3,4,4,4,4)$ & $0$ & $072155\ldots$\\
$20$ & $(4,4,4,4,4)$ & $0$ & $088189\ldots$\\
$21$ & $(4,4,4,4,5)$ & $0$ & $018452\ldots$\\
$22$ & $(4,4,4,5,5)$ & $0$ & $26319\ldots$\\
$23$ & $(4,5,4,5,5)$ & $0$ & $0040481\ldots$\\
$24$ & $(4,5,5,5,5)$ & $0$ & $060271\ldots$\\
$25$ & $(5,5,5,5,5)$ & $1$ & $7271\ldots\times 10^{-6}$\\
\bottomrule
\end{tabular}
\caption{Selected further counterexamples among blow-ups of $C_5$. The
displayed ratios are rounded from exact rational computations.}
\label{tab:c5more}
\end{table}

For the positive vectors in Table~\ref{tab:c5more}, the five parts are the
twin classes, so $\lvert\Aut(C_5[m])\rvert=(\prod_{i=1}^5m_i!)
\lvert\operatorname{Stab}_{\Aut(C_5)}(m)\rvert$; together with \eqref{eq:blowup}, this
gives every ratio in the table exactly.
Within the $C_5[m]$ family, $\Kbal$ is again smaller at $n=16$ and $n=18$,
where the five parts cannot all have equal size; the failure of the conjecture
is therefore not monotone in $n$ at this scale. The unrestricted cases
$n=16$ and $n=18$ remain open.

\section{Asymptotics}\label{sec:asym}

We first evaluate $\Lk(K_{m,m})$. Iterating \eqref{eq:brec} downwards from $(m,m)$ expresses $b^{(2m)}_m$ as a sum over lattice paths. Formally, call a \emph{deletion path} from $(a,b)$ a sequence of steps each decreasing one coordinate by $1$, terminating when a coordinate reaches $0$, and give the step that decreases the first coordinate at state $(a',b')$ the cost $\log_2\bin{a'+b'-1}{b'}$, the other step the cost $\log_2\bin{a'+b'-1}{a'}$. Unwinding \eqref{eq:brec} and using $b^{(k)}_0=1$,
\begin{equation}\label{eq:pathsum}
b^{(2m)}_m=\sum_{\pi}2^{-\mathrm{cost}(\pi)},
\end{equation}
the sum being over all deletion paths from $(m,m)$. Let $W(m)=\min_\pi \mathrm{cost}(\pi)$.

\begin{lemma}\label{lem:greedy}
$W(m)=\sum_{j=0}^{m-1}\log_2\bin{m+j}{j}$, attained by the path that empties one side first.
\end{lemma}

\begin{proof}
Let $f(a,b)$ denote the minimum cost of a deletion path from $(a,b)$, with $f(a,0)=f(0,b)=0$. We show by induction on $a+b$ that for $a\le b$,
\[
f(a,b)=S(a,b):=\sum_{j=0}^{a-1}\log_2\bin{b+j}{j},
\]
which is the cost of emptying the first side; by symmetry this suffices. The cases $a=0$ and $a+b\le2$ are clear. Let $1\le a\le b$. Emptying the first side has cost $\log_2\bin{a+b-1}{b}+S(a-1,b)=S(a,b)$, since $\bin{a+b-1}{b}=\bin{a+b-1}{a-1}$. So $f(a,b)\le S(a,b)$, and it remains to exclude the other first step, that is to show
\begin{equation}\label{eq:exchange}
S(a,b)\;\le\; \log_2\bin{a+b-1}{a}+f(a,b-1).
\end{equation}
Suppose first $a\le b-1$, so that $f(a,b-1)=S(a,b-1)$ by induction. Since
$\bin{b+j}{j}\big/\bin{b-1+j}{j}=(b+j)/b$, we get
\begin{align*}
S(a,b)-S(a,b-1)
&=\sum_{j=0}^{a-1}\log_2\frac{b+j}{b}\\
&=\log_2\frac{\prod_{j=0}^{a-1}(b+j)}{b^{a}},\\
\log_2\bin{a+b-1}{a}
&=\log_2\frac{\prod_{j=0}^{a-1}(b+j)}{a!},
\end{align*}
so \eqref{eq:exchange} is equivalent to $a!\le b^{a}$, which holds because $a\le b$. If instead $a=b$, then $f(a,b-1)=S(a-1,a)$ by induction and
$S(a,a)-S(a-1,a)=\log_2\bin{2a-1}{a-1}=\log_2\bin{2a-1}{a}$, so \eqref{eq:exchange} holds with equality, as symmetry demands.
\end{proof}

\begin{lemma}\label{lem:W}
$W(m)=\big(2-\tfrac{1}{2\ln 2}\big)m^2+O(m\log m)$.
\end{lemma}

\begin{proof}
By Lemma~\ref{lem:greedy}, $W(m)\ln 2=\sum_{j=0}^{m-1}\big(\ln (m+j)!-\ln m!-\ln j!\big)$. Stirling gives $\ln k!=k\ln k-k+O(\log k)$, and comparing the three sums with the integrals $\int x\ln x\,dx$ yields
\[
\sum_{j=0}^{m-1}\ln(m+j)!=\tfrac32m^2\ln m+(2\ln 2-\tfrac94)m^2+O(m\log m),
\]
\begin{align*}
m\ln m!&=m^2\ln m-m^2+O(m\log m),\\
\sum_{j=0}^{m-1}\ln j!&=\tfrac12m^2\ln m-\tfrac34m^2+O(m\log m).
\end{align*}
The $m^2\ln m$ terms cancel and the remainder is $(2\ln2-\tfrac12)m^2+O(m\log m)$.
\end{proof}

\begin{proof}[Proof of Theorem~\ref{thm:asym}]
There are at most $\bin{2m}{m}\le4^m$ deletion paths from $(m,m)$, so
\eqref{eq:pathsum} gives
\begin{align*}
2^{-W(m)}&\le b^{(2m)}_m\le 4^m2^{-W(m)},\\
\log_2(1/b^{(2m)}_m)&=W(m)+O(m).
\end{align*}
By \eqref{eq:Lb},
\[
\log_2\frac{1}{\Lk(K_{m,m})}
=\log_2\frac{1}{b^{(2m)}_m}+\log_2(2\times(2m)!).
\]
Lemma~\ref{lem:W} with $n=2m$ and $m^2=n^2/4$ therefore gives
\[
\log_2\frac1{\Lk(K_{m,m})}=\Big(\tfrac12-\tfrac1{8\ln2}\Big)n^2+O(n\log n).
\]
For odd $n$ the same argument applied to $(m,m+1)$ changes only the $O(n\log n)$ term.

For the second statement, let $g(n)$ be the number of isomorphism classes of graphs of order $n$. Since $\Lk$ is a probability distribution on these classes, $\min_G\Lk(G)\le 1/g(n)$, and $g(n)\ge 2^{\bin n2}/n!$, so $\log_2(1/\min_G\Lk)\ge\bin n2-\log_2n!=\tfrac{n^2}2+O(n\log n)$. Conversely \eqref{eq:BMSbounds} with $\lvert\Aut(G)\rvert\le n!$ gives $\min_G\Lk\ge 1/(n!P_n)$, and $\log_2P_n=\sum_{i=1}^n\log_2\bin{i-1}{\lfloor (i-1)/2\rfloor}=\bin n2+O(n\log n)$, whence $\log_2(1/\min_G\Lk)\le \tfrac{n^2}2+O(n\log n)$.

Finally $\tfrac12-\tfrac1{8\ln2}<\tfrac12$ strictly, so the two exponents differ in the leading term and the ratio is $2^{(1/(8\ln2)+o(1))n^2}$; in particular $\Lk(\Kbal)>\min_G\Lk(G)$ for all large $n$.
\end{proof}

\begin{remark}
Numerically $\log_2(1/\Lk(\Kbal))/(n^2/2)$ equals $0.8158$, $0.7515$, $0.7130$, $0.6852$ and $0.6554$ for $n=10,40,80,160,640$, descending towards $2\times(\tfrac12-\tfrac1{8\ln2})=1-\tfrac1{4\ln2}=0.6393\ldots$, and $W(m)/m^2$ equals $1.2615$ at $m=320$ against the limit $2-\tfrac1{2\ln 2}=1.27865\ldots$.
\end{remark}

\section{Entropy}\label{sec:entropy}

\begin{proof}[Proof of Theorem~\ref{thm:entropy}]
The random choices are independent across steps, so the entropy of the whole sequence of choices is
\[
\mathcal{H}_n=\sum_{t=1}^{n}\left(\log_2 t+
\mathbb{E}_k\log_2\bin{t-1}{k}\right),
\]
with $k$ uniform on $\{0,\dots,t-1\}$. Since
$\log_2\bin Nk=NH_2(k/N)+O(\log N)$ uniformly, where $H_2$ is the
binary entropy in bits, and $\int_0^1H_2(x)\,dx=1/(2\ln2)$, we get
\[
\mathbb{E}_k\log_2\bin{t-1}{k}=\frac{t}{2\ln2}+O(\log t),
\qquad
\mathcal H_n=\frac{n^2}{4\ln2}+O(n\log n).
\]
The isomorphism class is a function of the choices, so its entropy is at most
$\mathcal H_n$; conversely, conditioned on the class, the number of choice
sequences producing it is at most $n!$, so the conditional entropy is at most
$\log_2n!=O(n\log n)$. The comparison with
$\log_2g(n)=\tfrac{n^2}{2}+O(n\log n)$ is as in the proof of
Theorem~\ref{thm:asym}.

It remains to justify the assertion about a typical output. Let $Y_n$ denote
the complete choice history and put
\[
Z_n=-\log_2\Pr(Y_n)
=\sum_{t=1}^{n}\left(\log_2t+
\log_2\bin{t-1}{K_t}\right),
\]
where the $K_t$ are independent and uniform on $\{0,\dots,t-1\}$. The
summands are independent and bounded in absolute value by $O(t)$, so
$\operatorname{Var}(Z_n)=O(n^3)$. Chebyshev's inequality and the preceding
entropy computation give
\[
Z_n=\frac{n^2}{4\ln2}+o_{\mathrm p}(n^2).
\]
Now set $I_n=-\log_2\Lk(G_n)$. Since $G_n$ is a function of $Y_n$,
\[
D_n:=Z_n-I_n=-\log_2\Pr(Y_n\mid G_n)\ge0.
\]
Moreover $\mathbb E D_n=H(Y_n\mid G_n)\le\log_2n!=O(n\log n)$, because
each isomorphism class has at most $n!$ choice histories. Markov's inequality
therefore gives $D_n=o_{\mathrm p}(n^2)$, and hence
$I_n/n^2\xrightarrow{\mathrm p}1/(4\ln2)$.
\end{proof}

Exact values of the entropy for $n\le 9$, computed from all isomorphism
classes, are
\[
0,1,1.918,3.198,4.793,6.841,9.444,12.717,16.729
\]
bits, against
\[
\log_2g(n)=0,1,2,3.46,5.09,7.29,10.03,13.59,18.07.
\]
The ratio decreases slowly from $0.939$ at $n=6$ to $0.925$ at $n=9$,
consistent with the limiting ratio $1/(2\ln 2)=0.7213\ldots$ implied by
Theorem~\ref{thm:entropy}.  The concentration statement implies that there is
a sequence $\varepsilon_n\downarrow0$ such that the typical set
\[
\mathcal T_n=\left\{G:\left|-\frac{1}{n^2}\log_2\Lk(G)
-\frac{1}{4\ln2}\right|\le\varepsilon_n\right\}
\]
satisfies $\Pr(G_n\in\mathcal T_n)=1-o(1)$.  Bounding the likelihood of each
member of $\mathcal T_n$ above and below and summing shows that
$\lvert\mathcal T_n\rvert=2^{n^2/(4\ln2)+o(n^2)}$.  Thus the process
concentrates on a $2^{-(1/2-1/(4\ln2))n^2+o(n^2)}$ fraction of all
isomorphism classes while still spreading over $2^{\Theta(n^2)}$ of them,
which is the precise sense in which the uniform monkey is a poor but not
degenerate sampler.

\section{The maximum, and open problems}\label{sec:remarks}

For the maximum no analogue of Theorem~\ref{thm:bip} is available, and the situation is genuinely irregular. Up to complementation, the maximisers of $\Lk$ are threshold graphs for $n\le 6$; for $n=7$ and $n=8$ they are split but not threshold; and for $n=9$ the maximiser is the disjoint union of two isolated vertices with a caterpillar on seven vertices, which is not split, being obtained from $P_6$ by attaching a pendant vertex at the third vertex. Thus neither the threshold graphs nor the broader class of split graphs contains the maximisers for all orders. The next three exact maximum values are
\[
\max_{\lvert V(G)\rvert=10}\Lk(G)=\frac{219903017}{4938071040000},
\]
\[
\max_{\lvert V(G)\rvert=11}\Lk(G)=\frac{7472883629}{1140694410240000},
\]
\[
\max_{\lvert V(G)\rvert=12}\Lk(G)=\frac{18096017946691}{17566693917696000000}.
\]
The complementary maximising pairs have graph6 representations
\verb|I????GICo| and \verb|Ihuz~~~~w| for $n=10$;
\verb|J?????A@OR?| and \verb|Jhuz~~~~~~_| for $n=11$; and
\verb+K????C?@?PAE+ and \verb+Kh~u||~~~~~~+ for $n=12$.
What does appear stable is the order of magnitude: $-\log_2\max_G\Lk(G)$ equals
\[
0,1,1.585,2.469,3.815,5.487,7.439,9.684,12.021,14.455,17.220,19.889
\]
for $n\le 12$, consistent with $\max_G\Lk(G)=2^{\Theta(n)}/n!$.

\begin{problem}\label{prob:min}
Determine $\arg\min_G\Lk(G)$ for $n\ge15$. In particular, is
$C_5[\overline{K_3}]$ a minimiser at $n=15$? Are the minimisers always
blow-ups, and can their base graphs be chosen from a fixed finite family, or
must their orders grow with $n$? Theorem~\ref{thm:asym} shows that
$\lvert\Aut\rvert$ is of secondary importance on the $n^2$ logarithmic scale
and suggests looking for graphs whose typical vertex orderings have back
degrees near the centre of their range.
\end{problem}

\begin{problem}
Settle Conjecture~\ref{conj:DMS} for the remaining small orders $n=16$ and
$n=18$.
\end{problem}

\begin{problem}
Prove $A(G)\le C^n\lvert\Aut(G)\rvert$ for an absolute constant $C$. Equivalently,
prove
\[
\max_G\Lk(G)=2^{O(n)}/n!.
\]
\end{problem}

\begin{problem}\label{prob:complexity}
Settle the complexity of $\Lk$, raised in \cite{BMS14} and repeated in \cite{DMS18}. Is exact evaluation of $A$ $\#\mathrm P$-hard? Corollary~\ref{cor:algo} is a dynamic program over vertex subsets, which suggests that $A$ is computable in polynomial time on graphs of bounded pathwidth.
\end{problem}

\begin{problem}
Study graph likelihood for a general attachment kernel, replacing the uniform
choice of $k$ by a distribution $\nu_t$ on $\{0,\dots,t-1\}$ as in
\cite{JS13}.  Assuming $\nu_1(0)=1$, the subset dynamic program extends by
putting $F_\nu(\{v\})=1$ and
\[
F_\nu(S)=\lvert S\rvert\sum_{v\in S}
\frac{\nu_{\lvert S\rvert}(\deg_{G[S]}v)F_\nu(S\setminus\{v\})}
{\bin{\lvert S\rvert-1}{\deg_{G[S]}v}},
\qquad
\Lk_\nu(G)=\frac{F_\nu(V(G))}{n!\,\lvert\Aut(G)\rvert}.
\]
Determine how the extrema and entropy depend on $\nu$.  The same recurrence
turns $\Lk_\nu$ into an exactly computable marginal likelihood of an
unlabelled graph, which is the quantity that likelihood-based inference for
mechanistic network models must otherwise approximate.
\end{problem}

\subsection*{Computational verification}

All computations use exact integer or rational arithmetic, with \texttt{geng}
and \texttt{nauty} \cite{nauty} for graph generation and automorphism groups.
The primary exhaustive certificate for Theorem~\ref{thm:small} is an
independently developed C plugin compiled directly into \texttt{geng}.  It
uses the exact automorphism-group order returned by \texttt{nauty}, checked
unsigned $128$-bit arithmetic for the integer-scaled subset recurrence, and
exact $192$-bit cross-products for rational likelihood comparisons.  It
rejects orders above $12$ rather than risk overflow.  The same program
reproduced the previously known extrema through order $10$ and independently
reobtained the order-$11$ result before the order-$12$ production run; thus the
claims through $n=12$ do not depend on
the preliminary screening programs supplied with the initial manuscript.

At order $12$, complement symmetry reduces the computation to one
representative from each of the $82\,545\,586\,656$ complement orbits of the
$165\,091\,172\,592$ unlabelled graphs; the latter count includes $720$
self-complementary graphs.  The computation was divided into $4\,800$ exact
\texttt{geng} residue shards and run with $14$ workers.  All shard outputs
completed, with no failed shards or nonempty error files.  Their equivalent
full graph counts sum to $165\,091\,172\,592$ before the merge reports the
global extrema.  Aggregate helper CPU time was about $817$ core-hours, and
elapsed wall time was approximately $68$ hours.

The principal implementation for individual larger graphs performs graph
processing in the Wolfram Language and invokes a C helper that computes the
scaled recurrence modulo multiple $64$-bit primes; the exact result is
reconstructed by the Chinese remainder theorem.  It returns
\[
A(C_5[\overline{K_3}])
=\frac{63977511069907}{427822618422142944000000},
\]
agreeing with \eqref{eq:blowup} evaluated over the $4^5$ vectors
$m'\le(3,3,3,3,3)$ and with the direct rational recurrence over the $2^{15}$
subsets.

The blow-up searches in Proposition~\ref{prop:first15} use a separate C
program that reads the base graphs streamed by \texttt{geng}, enumerates the
specified positive part vectors, and evaluates an integer-scaled form of
\eqref{eq:blowup}.  Dynamic-programming quantities use checked unsigned
$256$-bit arithmetic, automorphism orders are computed exactly by
\texttt{nauty}, and likelihood comparisons use exact unsigned $320$-bit
cross-products.  The program exits rather than returning a value if any
checked bound is exceeded.  It reproduces the likelihoods of the known
minimising graphs through $n=12$, the value of
$C_5[\overline{K_3}]$ above, and the $n=18$ rook-graph ratio reported after
Proposition~\ref{prop:first15}.

For the extended $C_5[m]$ calculation, a separate arbitrary-precision
implementation enumerates positive five-part vectors modulo the dihedral
symmetries of $C_5$ and evaluates \eqref{eq:blowup} with exact rationals.  It
reproduces the ratios in Table~\ref{tab:c5more}, exhaustively checks every
positive vector through order $50$, and evaluates every nearly balanced
vector through order $100$.  Independent exact evaluations in the Wolfram
Language reproduce the resulting ratios at orders $26$, $30$ and $50$.

For each base order $p$ and target order $n$, the search checks the
$\binom{n-1}{p-1}$ positive ordered part vectors for every base.  The rows
through $p=9$ are exhaustive, while those for $10\le p\le12$ apply the exact
necessary screens described above before likelihood evaluation.  The
screened order-$11$ and order-$12$ bases from the $n=13$ searches are reused
for $n=14$.  The order-$12$ stream was divided into $4\,096$ shards; the
merged input counts sum to all $165\,091\,172\,592$ unlabelled bases and the
retained file contains $140\,474$ distinct graph6 strings.  The order-$12$,
$n=13$ CPU entry includes $125\,939.981$ seconds for the shared base screen
and $67.805$ seconds for the composition search.
Parallel runs are merged only after verifying the total numbers of bases and
blow-ups.  Table~\ref{tab:blowupsearch} records
all exact composition searches used here.  CPU times are aggregate
process times for the blow-up search program; they exclude CPU time used by
\texttt{geng} for graph generation and are not elapsed wall times.

\begin{table}[htbp]
\centering
\small
\begin{tabular}{@{}rrrrrr@{}}
\toprule
$n$ & $p$ & bases & vectors/base & likelihood evals & CPU seconds\\
\midrule
$13$ & $7$ & $1\,044$ & $924$ & $964\,656$ & $13.685$\\
$14$ & $7$ & $1\,044$ & $1\,716$ & $1\,791\,504$ & $33.791$\\
$16$ & $7$ & $1\,044$ & $5\,005$ & $5\,225\,220$ & $171.449$\\
$18$ & $7$ & $1\,044$ & $12\,376$ & $12\,920\,544$ & $625.523$\\
$13$ & $8$ & $12\,346$ & $792$ & $9\,778\,032$ & $186.062$\\
$14$ & $8$ & $12\,346$ & $1\,716$ & $21\,185\,736$ & $557.457$\\
$16$ & $8$ & $12\,346$ & $6\,435$ & $34\,372\,082$ & $1\,440.603$\\
$18$ & $8$ & $12\,346$ & $19\,448$ & $156\,498\,056$ & $12\,447.239$\\
$13$ & $9$ & $274\,668$ & $495$ & $135\,960\,660$ & $4\,078.112$\\
$14$ & $9$ & $274\,668$ & $1\,287$ & $353\,497\,716$ & $17\,078.575$\\
$13$ & $10$ & $12\,005\,168$ & $220$ & $528\,765$ & $31.630$\\
$14$ & $10$ & $12\,005\,168$ & $715$ & $80\,258$ & $13.454$\\
$13$ & $11$ & $1\,018\,997\,864$ & $66$ & $555\,612$ & $598.681$\\
$14$ & $11$ & $1\,018\,997\,864$ & $286$ & $79\,908$ & $9.805$\\
$13$ & $12$ & $165\,091\,172\,592$ & $12$ & $758\,030$ & $126\,007.786$\\
$14$ & $12$ & $165\,091\,172\,592$ & $78$ & $121\,413$ & $18.206$\\
\bottomrule
\end{tabular}
\caption{Exact independent-blow-up searches.  Here $p$ is the base order;
the order-$8$ rows for $n=13,14$ record intermediate staged checks.  Rows through $p=9$
evaluate every base/part-vector representation.  The rows with
$10\le p\le12$ use the necessary automorphism screens described in the text;
their bases and vectors/base columns give the full unscreened mathematical
input counts.  The $p=11$ and $p=12$,
$n=14$ searches reused the bases saved during the corresponding $n=13$
streams, so their CPU times exclude the shared initial scans.}
\label{tab:blowupsearch}
\end{table}

The remaining order-$13$ and order-$14$ searches use the complete tables of
marks for $\mathfrak S_{12}$ and $\mathfrak S_{13}$ in \textsc{TomLib}.  For
$\mathfrak S_{14}$, subgroup classes above the required threshold are obtained
by recursively taking maximal-subgroup class representatives and identifying
subgroups conjugate in $\mathfrak S_{14}$.  The traversal has $6280$ classes
and no failed maximal-subgroup computations.  A separate C program expands
the group orbits on unordered pairs using disjoint checked unsigned $128$-bit
masks and applies the exact twin screens described above.  Graphs are
canonically labelled by \texttt{nauty}, deduplicated, and passed to the same
checked likelihood program used for the blow-up computations.  Independent
residue shards sum to the stated orbital representation counts before the
canonical merge.

For the partial order-$16$ calculation, a checkpointed maximal-subgroup
traversal above automorphism threshold $1556$ produces $116\,597$ subgroup
classes and no failed maximal-subgroup computations.  The forced-twin test
leaves $44$ twin-free-possible classes and $4580$ orbital representations.
Direct twin testing leaves $738$ representations, canonical labelling leaves
$128$ graphs, and exact likelihood evaluation finds no graph below
$K_{8,8}$.  This certifies the twin-free part only; the $116\,553$ subgroup
classes that force a twin pair are not expanded into graphs by this
calculation.

For the vertex-transitive search, the GAP 4.15.1 transitive group library was
used \cite{GAP}.  At orders $13$ and $14$ there
are $9$ and $63$ transitive permutation groups.  Enumerating all unions of
their orbits on unordered vertex pairs gives $152$ and $1684$ graph
representations, which canonical labelling reduces to $14$ and $56$
isomorphism classes.  Likelihoods are then evaluated by the integer-scaled
subset recurrence above.

The degree-sequence calculation enumerates all graphical sequences and applies
Lemma~\ref{lem:screen} by integer cross multiplication.
For its second deletion, the neighbours of each possible deleted vertex are
assigned among the remaining degree classes so as to minimise the resulting
lower bound; allowing the minimizing assignment independently for each
vertex is a relaxation and therefore cannot discard a genuine competitor.

\begin{thebibliography}{99}
\small
\raggedright
\setlength{\itemsep}{0pt}

\bibitem{BMS14} C.~R.~S. Banerji, T. Mansour and S. Severini, \emph{A notion of graph likelihood and an infinite monkey theorem}, arXiv:1304.3600 (2013); J.\ Phys.\ A: Math.\ Theor.\ \textbf{47} (2014) 035101, \url{https://doi.org/10.1088/1751-8113/47/3/035101}.

\bibitem{DMS18} D. Dervovic, A. Mocherla and S. Severini, \emph{Constructing graphs with limited resources}, arXiv:1802.09844 (2018), \url{https://arxiv.org/abs/1802.09844}.

\bibitem{JS13} S. Janson and S. Severini, \emph{An example of graph limits of growing sequences of random graphs}, J.\ Comb.\ \textbf{4} (2013), no.~1, 67--80, \url{https://doi.org/10.4310/JOC.2013.v4.n1.a3}.

\bibitem{nauty} B.~D. McKay and A. Piperno, \emph{Practical graph isomorphism, II}, J.\ Symbolic Comput.\ \textbf{60} (2014) 94--112, \url{https://doi.org/10.1016/j.jsc.2013.09.003}.

\bibitem{OEIS} OEIS Foundation Inc., \emph{The On-Line Encyclopedia of Integer Sequences}, \url{https://oeis.org}; minimum and maximum likelihood by vertex count: \href{https://oeis.org/A234234}{A234234}/\href{https://oeis.org/A234235}{A234235} and \href{https://oeis.org/A234236}{A234236}/\href{https://oeis.org/A234237}{A234237}, respectively.

\bibitem{GK22} M. Grech and A. Kisielewicz, \emph{Wreath product in automorphism groups of graphs}, J.\ Graph Theory \textbf{101} (2022) 29--51, \url{https://doi.org/10.1002/jgt.22808}.

\bibitem{GAP} The GAP Group, \emph{GAP---Groups, Algorithms, and Programming}, Version 4.15.1 (2025), \url{https://www.gap-system.org}.

\bibitem{TomLib} T. Merkwitz, L. Naughton and G. Pfeiffer, \emph{TomLib---The GAP Library of Tables of Marks}, Version 1.2.11 (2024), \url{https://gap-packages.github.io/tomlib}.

\bibitem{Sab59} G. Sabidussi, \emph{The composition of graphs}, Duke Math.\ J.\ \textbf{26} (1959) 693--696, \url{https://doi.org/10.1215/S0012-7094-59-02667-5}.

\bibitem{Sab61} G. Sabidussi, \emph{The lexicographic product of graphs}, Duke Math.\ J.\ \textbf{28} (1961) 573--578, \url{https://doi.org/10.1215/S0012-7094-61-02857-5}.

\bibitem{MW} E.~W. Weisstein, \emph{Graph likelihood}, MathWorld---A Wolfram Resource, \url{https://mathworld.wolfram.com/GraphLikelihood.html}.

\bibitem{GraphData} Wolfram Research, \texttt{GraphData}, Wolfram Language
documentation, \url{https://reference.wolfram.com/language/ref/GraphData.html}.

\end{thebibliography}

\end{document}
